\chapter{Stone--von Neumann 定理}
\label{chap:stone-von-neumann}
\index{Stone--von Neumann 定理}
\index{Segal--Bargmann 空间}

Stone--von Neumann 定理是关于满足正则对易关系的算子的唯一性定理.设 $A,B$ 是 $\mathcal H$ 上满足 $[A,B]=i\hbar I$ 的两个自伴算子,并且 $A,B$ 在 $\mathcal H$ 上不可约作用,即只有 $\{0\}$ 与 $\mathcal H$ 是同时在 $A,B$ 下不变的闭子空间.在适当的技术假设(指数化对易关系)下,我们将证明 $A,B$ 与通常的位置、动量算子 $X,P$ 酉等价:存在酉算子 $U:\mathcal H\to L^2(\mathbb R)$,使 $UAU^{-1}=X$、$UBU^{-1}=P$.如果 $\mathcal H$ 上的作用不是不可约的,则 $\mathcal H$ 分解为 $A,B$ 的不变子空间 $V_l$ 的直和,并且 $A,B$ 在每个 $V_l$ 上的限制都分别与通常的 $X,P$ 酉等价.

我们先给出 Stone--von Neumann 定理的启发式论证,其中略去一些本质但技术性的定义域问题;然后引入可以看作普通正则对易关系的温和加强的指数化对易关系;最后给出定理的精确表述与证明.

\section{启发式论证}
\label{sec:stone-heuristic}

设 $A,B$ 是可分 Hilbert 空间 $\mathcal H$ 上满足 $[A,B]=i\hbar I$ 的两个(可能无界的)自伴算子.我们希望证明,$\mathcal H$ 看起来像由一族在 $A,B$ 下不变的闭子空间 $V_l$ 构成的 Hilbert 空间直和,并且每个 $V_l$ 都与 $L^2(\mathbb R)$ 酉等价,酉等价把 $A,B$ 变为标准位置与动量算子.也就是说,希望找到酉映射 $U_l:V_l\to L^2(\mathbb R)$,使
\[
U_lAU_l^{-1}=X,
\qquad U_lBU_l^{-1}=P.
\]

由于相关无界算子的定义域问题,这一结论并不完全正确.不过可以先作如下启发式论证.仿照第~\ref{chap:harmonic-oscillator} 章中 $a,a^*$ 的定义,取升降算子
\[
\alpha=\frac{m\omega A+iB}{\sqrt{2\hbar m\omega}},
\qquad
\alpha^*=\frac{m\omega A-iB}{\sqrt{2\hbar m\omega}}.
\]
考虑降算子 $\alpha$ 的核 $W$;只要 $\alpha$ 是闭算子,$W$ 就是闭子空间,可以把其中元素看作 $\alpha^*\alpha$ 的“基态”.取 $W$ 的正交归一基 $\{\phi_0^l\}$,定义
\[
\phi_m^l=(\alpha^*)^m\phi_0^l.
\]
不难看出,当 $l\neq l'$ 时,$\phi_m^l$ 与 $\phi_{m'}^{l'}$ 对所有 $m,m'$ 都正交.令 $V_l$ 为固定 $l$ 时所有 $\phi_m^l$ 的闭线性张成.

由第~\ref{sec:harmonic-algebraic} 节的计算,$\alpha,\alpha^*$ 在每条链 $\{\phi_m^l\}_{m\geq0}$ 上的作用方式,与标准降、升算子 $a,a^*$ 在 $a^*a$ 的特征向量链上的作用方式完全相同.因此,可以定义从 $V_l$ 到 $L^2(\mathbb R)$ 的酉映射 $U_l$,把 $\phi_m^l$ 映到定理~\ref{thm:harmonic-completeness} 中的谐振子特征函数 $\psi_m$.于是 $U_l$ 把 $\alpha,\alpha^*$ 分别交织到 $a,a^*$,也把 $A,B$ 交织到 $X,P$.

还需(启发式地)说明所有 $V_l$ 填满 $\mathcal H$.令 $V$ 为这些子空间的张成;$V$ 在 $\alpha,\alpha^*$ 下不变,所以 $V^\perp$ 在伴随算子 $\alpha^*,\alpha$ 下也不变.若 $V^\perp\neq\{0\}$,仿照第 11 章的论证,$V^\perp$ 中应存在一个被 $\alpha$ 消去的非零基态.但这会与 $\{\phi_0^l\}$ 构成 $\ker\alpha$ 的正交归一基相矛盾.

上述论证不可能完全严格,因为第~\ref{sec:uncertainty-counterexample} 节的反例给出了一对满足正则对易关系、但显然不可能与通常位置、动量算子酉等价的算子:该节中的“位置”算子 $A$ 是有界的,而通常的位置算子是无界的.问题仍然来自定义域.令 $m,\hbar,\omega$ 都为 $1$,寻找被
\[
\alpha=\frac1{\sqrt2}\left(A+iB\right)=\frac1{\sqrt2}\left(x+\frac{d}{dx}\right)
\]
消去的向量.与第 11 章相同的论证说明,$\phi_0$ 必须是 $e^{-x^2/2}$ 的常数倍;于是 $\phi_1:=\alpha^*\phi_0$ 是 $xe^{-x^2/2}$ 的倍数.但 $\phi_1$ 不属于 $\alpha^*$ 的定义域,因为它不满足定义 $B$ 时的周期边界条件 $\psi(-1)=\psi(1)$.所以无法继续施加 $\alpha^*$,正交链的构造完全崩溃.

因此需要某种额外条件,把正则对易关系的“好”情形与“坏”情形区分开来.一种可能的条件是正则对易关系的指数化形式,下一节将讨论它.本章第~\ref{sec:stone-theorem} 节的严格证明沿用上述启发式论证的结构,只是用通过 Weyl 量子化类似构造的有界算子替代无界的 $\alpha,\alpha^*$.

\section{指数化对易关系}
\label{sec:exponentiated-commutation}

若 $A$ 是 Hilbert 空间 $\mathcal H$ 上的有界算子,可以用幂级数定义指数
\[
e^A=\sum_{m=0}^{\infty}\frac{A^m}{m!},\qquad A^0=I.
\]
标准的幂级数论证表明,若 $A,B\in\mathcal B(\mathcal H)$ 对易,则
\[
e^{A+B}=e^Ae^B,qquad [A,B]=0.
\label{eq:exp-commuting}
\]
即使 $A,B$ 不对易,也有 Baker--Campbell--Hausdorff 公式;当 $A,B$ 足够小时,$e^Ae^B$ 可表示为
\[
e^Ae^B=\exp\left(A+B+\frac12[A,B]+\frac1{12}[A,[A,B]]+\cdots\right).
\]
当 $A,B$ 与其对易子都对易时,所有更高阶对易子消失.

\begin{theorem}[特殊的 Baker--Campbell--Hausdorff 公式]\label{thm:bch-special}
设 $A,B\in\mathcal B(\mathcal H)$,且 $A,B$ 与其对易子对易,即
\[
[A,[A,B]]=[B,[A,B]]=0.
\]
则
\[
e^Ae^B=e^{A+B+[A,B]/2},
\qquad
e^{A+B}=e^{-[A,B]/2}e^Ae^B.
\]
在这一特殊情形中,不需要对 $A,B$ 施加小性假设.
\end{theorem}

\begin{proof}
证明
\[
e^{tA}e^{tB}=e^{t(A+B)+t^2[A,B]/2}
\label{eq:bch-t}
\]
即可在 $t=1$ 时得到结论.由于 $[A,B]$ 与所有相关算子对易,式~\eqref{eq:bch-t} 等价于
\[
e^{tA}e^{tB}e^{-t^2[A,B]/2}=e^{t(A+B)}.
\label{eq:bch-equivalent}
\]
令左端为 $\alpha(t)$.对任意有界 $C$,逐项求导可得
\[
\frac{d}{dt}e^{tC}=Ce^{tC}=e^{tC}C,
\qquad
\frac{d}{dt}e^{-t^2[A,B]/2}=e^{-t^2[A,B]/2}(-t[A,B]).
\]
由乘积法则,
\[
\frac{d\alpha}{dt}=e^{tA}Ae^{tB}e^{-t^2[A,B]/2}
+e^{tA}e^{tB}Be^{-t^2[A,B]/2}
+e^{tA}e^{tB}e^{-t^2[A,B]/2}(-t[A,B]).
\]
还需一个中间公式.由乘积法则与 $B$ 和 $[A,B]$ 对易,
\[
\frac{d}{dt}\left(e^{-tB}Ae^{tB}\right)=e^{-tB}[A,B]e^{tB}=[A,B].
\]
在 $t=0$ 处取值为 $A$,积分得
\[
e^{-tB}Ae^{tB}=A+t[A,B].
\label{eq:bch-conjugation}
\]
于是
\[
e^{tA}Ae^{tB}=e^{tA}e^{tB}(A+t[A,B]).
\]
把这一关系代回,并利用 $[A,B]$ 与所有因子对易,得到
\[
\frac{d\alpha}{dt}=\alpha(t)(A+t[A,B]+B-t[A,B])=\alpha(t)(A+B).
\]
由练习~\ref{ex:stone-ode},该方程的唯一解为 $\alpha(t)=\alpha(0)e^{t(A+B)}$;而 $\alpha(0)=I$,故式~\eqref{eq:bch-equivalent} 成立.
\end{proof}

\def\ChapterFourteenSegalLater{
\subsection{指数化对易关系}
\label{subsec:segal-exponentiated}

在 Segal--Bargmann 空间上定义自伴的“位置”与“动量”算子
\[
A_j=\frac1{\sqrt2}\left(z_j+\hbar\frac{\partial}{\partial z_j}\right),
\qquad
B_j=\frac{i}{\sqrt2}\left(z_j-\hbar\frac{\partial}{\partial z_j}\right).
\label{eq:segal-position-momentum}
\]
我们将找出以这些算子(的延拓)为无穷小生成元的一参数酉群;由 Stone 定理,生成元在适当定义域上确实自伴.然后验证指数化对易关系与不可约性.

为启发式地计算,形式地把定理~\ref{thm:bch-special} 应用于无界算子 $-\bar a_jz_j$ 与 $\hbar a_j\partial/\partial z_j$,得到
\[
\exp\left(\sum_{j=1}^n\left(-\bar a_jz_j+\hbar a_j\frac{\partial}{\partial z_j}\right)\right)
=e^{-\hbar|a|^2/2}
\exp\left(-\sum_{j=1}^n\bar a_jz_j\right)
\exp\left(\hbar\sum_{j=1}^na_j\frac{\partial}{\partial z_j}\right).
\label{eq:segal-formal-exponential}
\]
这提示定义
\[
(T_aF)(z)=e^{-\hbar|a|^2/2}e^{-\bar a\cdot z}F(z+\hbar a),
\qquad a\in\mathbb C^n,
\label{eq:segal-translation}
\]
其中对 $a,b\in\mathbb C^n$,约定 $a\cdot b=\sum_ja_jb_j$,不取复共轭.式~\eqref{eq:segal-formal-exponential} 左端是斜自伴的(算子与其伴随之差),所以预期 $T_a$ 是酉算子;对 $a$ 取适当实方向,左端就给出 $A_j$ 或 $B_j$ 生成的一参数群.

\begin{theorem}[Segal--Bargmann 平移算子]\label{thm:segal-translation}
对每个 $a\in\mathbb C^n$,式~\eqref{eq:segal-translation} 定义的 $T_a$ 是 Segal--Bargmann 空间上的酉算子,且 $a\mapsto T_a$ 强连续.它们满足
\[
T_aT_b=e^{i\hbar\operatorname{Im}(\bar a\cdot b)}T_{a+b}.
\label{eq:segal-translation-product}
\]
特别地,
\[
U_j(t):=T_{it e_j/\sqrt2},
\qquad
V_j(t):=T_{t e_j/\sqrt2}
\]
是强连续的一参数酉群;其无穷小生成元分别为 $A_j,B_j$,并满足指数化对易关系.

若 $F\in\operatorname{Dom}(A_j)$,则
\[
(A_jF)(z)=\frac1{\sqrt2}\left(z_jF(z)+\hbar\frac{\partial F}{\partial z_j}(z)\right),
\]
若 $F\in\operatorname{Dom}(B_j)$,则
\[
(B_jF)(z)=\frac{i}{\sqrt2}\left(z_jF(z)-\hbar\frac{\partial F}{\partial z_j}(z)\right).
\]
两者的定义域都包含所有全纯多项式.最后,$A_j,B_j$ 在 Segal--Bargmann 空间上不可约作用.
\end{theorem}

\begin{proof}
对固定 $a$,$T_aF$ 显然关于 $z$ 全纯.若 $F\in\mathcal H L^2(\mathbb C^n,\mu_\hbar)$,则
\[
\begin{aligned}
\|T_aF\|_\hbar^2
&=(\pi\hbar)^{-n}\int_{\mathbb C^n}e^{-\hbar|a|^2-2\operatorname{Re}(\bar a\cdot z)}
|F(z+\hbar a)|^2e^{-|z|^2/\hbar}\,dz\\
&=(\pi\hbar)^{-n}\int_{\mathbb C^n}|F(z+\hbar a)|^2e^{-|z+\hbar a|^2/\hbar}\,dz
=\|F\|_\hbar^2.
\end{aligned}
\]
所以 $T_a$ 是等距的.直接计算得到式~\eqref{eq:segal-translation-product}(见练习~\ref{ex:segal-translation-product});由此 $T_aT_{-a}=I$,故 $T_a$ 满射,从而酉.强连续性可先在全纯多项式上用支配收敛定理验证,再由多项式的稠密性推广到整个空间.

式~\eqref{eq:segal-translation-product} 立即给出 $U_j,V_j$ 是一参数酉群,并且这些群的生成元满足指数化对易关系.若 $F$ 属于 $U_j$ 生成元的定义域,则
\[
(A_jF)(z)=\lim_{t\to0}\frac1{it}\left(e^{-\hbar t^2/4}e^{itz_j/\sqrt2}F(z+it\hbar e_j/\sqrt2)-F(z)\right).
\]
该 $L^2$ 极限与逐点极限一致,得到式中 $A_j$ 的公式;全纯多项式上可用支配收敛定理验证极限确实存在.$B_j$ 同理.

最后证明不可约性.若 $A_j,B_j$ 不可约,则将定理~\ref{thm:stone-von-neumann} 应用于本空间时会得到至少两个不变子空间,从而得到至少两个线性无关的全纯函数 $F_l$,它们属于所有相关定义域且满足
\[
0=(A_j+iB_j)F_l=\sqrt2\hbar\frac{\partial F_l}{\partial z_j}
\]
对每个 $j$ 成立.这意味着每个 $F_l$ 都是常数,与线性无关性矛盾.
\end{proof}

\subsection{再生核}
\label{subsec:segal-reproducing-kernel}

由式~\eqref{eq:segal-point-evaluation},在固定点 $z$ 处取值是 $\mathcal H L^2(\mathbb C^n,\mu_\hbar)$ 上的连续线性泛函.因此它可以写成与唯一元素 $\chi_z$ 的内积.$\chi_z$ 称为参数为 $z$ 的相干态.

\begin{proposition}[Segal--Bargmann 再生核]\label{prop:segal-reproducing-kernel}
对任意 $F\in\mathcal H L^2(\mathbb C^n,\mu_\hbar)$,有
\[
F(z)=\int_{\mathbb C^n}e^{z\cdot\bar w/\hbar}F(w)\mu_\hbar(w)\,dw.
\label{eq:segal-reproducing}
\]
核 $e^{z\cdot\bar w/\hbar}$ 称为再生核,因为对它积分会“再生”函数 $F$.等价地,
\[
F(z)=\langle\chi_z,F\rangle_\hbar,
\qquad
\chi_z(w)=e^{\bar z\cdot w/\hbar}.
\]
\end{proposition}

\begin{proof}
先取 $z=0$.在命题~\ref{prop:segal-bargmann-complete} 的证明中,$F$ 的 Taylor 级数在该 Hilbert 空间中收敛,且不同单项式正交.因此 $\langle1,F\rangle_\hbar$ 只保留常数项;由于 $\mu_\hbar$ 是概率测度,
\[
\langle1,F\rangle_\hbar=F(0).
\]
对 $T_aF$ 应用同一关系.由定理~\ref{thm:segal-translation},$T_a^{-1}=T_{-a}$,所以
\[
(T_aF)(0)=\langle1,T_aF\rangle_\hbar=\langle T_{-a}1,F\rangle_\hbar.
\]
写出积分并令 $a=z/\hbar$,整理即得式~\eqref{eq:segal-reproducing}.
\end{proof}

\subsection{Segal--Bargmann 变换}
\label{subsec:segal-bargmann-transform}

定理~\ref{thm:segal-translation} 中的 $A_j,B_j$ 满足指数化对易关系并在 Segal--Bargmann 空间上不可约作用.因此 Stone--von Neumann 定理的第二部分给出一个除整体相位外唯一的酉映射
\[
U:\mathcal H L^2(\mathbb C^n,\mu_\hbar)\to L^2(\mathbb R^n),
\]
把这些算子交织到通常的位置、动量算子.其逆映射 $V:L^2(\mathbb R^n)\to\mathcal H L^2(\mathbb C^n,\mu_\hbar)$ 称为 Segal--Bargmann 变换.

\begin{theorem}[Segal--Bargmann 变换公式]\label{thm:segal-bargmann-transform}
令 $V$ 为上述酉映射的逆,并归一化使得 $V$ 把式~\eqref{eq:stone-ground-state} 中取 $\sigma=\hbar$ 的基态 $\phi_0$ 映到常函数 $1$.则
\[
(V\psi)(z)=(\pi\hbar)^{-n/4}\int_{\mathbb R^n}
\exp\left(-\frac1{2\hbar}\bigl(z\cdot z-2\sqrt2z\cdot x+x\cdot x\bigr)\right)\psi(x)\,dx.
\label{eq:segal-bargmann-transform}
\]
这里 $a\cdot b=\sum_ja_jb_j$ 不取复共轭,所以对每个固定 $x$,被积函数关于 $z$ 全纯.
\end{theorem}

当 $z=0$ 时,$(V\psi)(0)$ 正是 $\psi$ 与 $\sigma=\hbar$ 时基态的内积.一般 $z$ 处的值是 $\psi$ 与基态的相空间平移的内积乘以一个常数;这些平移态都是最小不确定态.

\begin{proof}
由 $V$ 的酉性与命题~\ref{prop:segal-reproducing-kernel} 在 $z=0$ 时的结论,
\[
\langle\phi_0,\psi\rangle_{L^2(\mathbb R^n)}=(V\psi)(0).
\]
更一般地,对 $a,b\in\mathbb R^n$,
\[
\begin{aligned}
\langle e^{-ia\cdot X}e^{-ib\cdot P}\phi_0,\psi\rangle
&=\langle\phi_0,e^{ib\cdot P}e^{ia\cdot X}\psi\rangle\\
&=\langle1,e^{ib\cdot B}e^{ia\cdot A}V\psi\rangle_\hbar
=(e^{ib\cdot B}e^{ia\cdot A}V\psi)(0).
\end{aligned}
\label{eq:segal-transform-intermediate}
\]
由式~\eqref{eq:segal-translation-product},右端等于
\[
e^{i\hbar a\cdot b/2}e^{-\hbar(|a|^2+|b|^2)/4}(V\psi)(\hbar(b+ia)/\sqrt2).
\]
取 $a=\sqrt2y_0/\hbar$、$b=\sqrt2x_0/\hbar$,得到
\[
\left\langle e^{-i\sqrt2y_0\cdot X/\hbar}e^{-i\sqrt2x_0\cdot P/\hbar}\phi_0,\psi\right\rangle
=e^{ix_0\cdot y_0/\hbar}e^{-(|x_0|^2+|y_0|^2)/(2\hbar)}
(V\psi)(x_0+iy_0).
\]
将左端写成 $L^2(\mathbb R^n)$ 上的积分并整理指数,正好得到式~\eqref{eq:segal-bargmann-transform}.
\end{proof}

\section{练习}
\label{sec:chapter14-exercises}

\begin{exercise}[半指数化对易关系]\label{ex:stone-semi-exponentiated}
证明:若算子 $A,B$ 满足第~\ref{sec:exponentiated-commutation} 节的指数化对易关系,则它们满足半指数化对易关系,即定理~\ref{thm:uncertainty-semi-exponentiated} 的假设.

提示:对 $a,s\in\mathbb R$ 与 $\psi\in\operatorname{Dom}(A)$,利用指数化对易关系整理
\[
\frac{e^{isA}(e^{iaB}\psi)-e^{iaB}\psi}{s},
\]
然后令 $s\to0$ 并使用 Stone 定理.
\end{exercise}

\begin{exercise}[有界算子微分方程]\label{ex:stone-ode}
\begin{enumerate}
\item 设 $\alpha:\mathbb R\to\mathcal B(\mathcal H)$ 可微,即每个 $t$ 处
\[
\lim_{h\to0}\frac{\alpha(t+h)-\alpha(t)}h
\]
在算子范数拓扑中存在.若 $d\alpha/dt=0$,证明 $\alpha$ 为常值.
\item 若 $\alpha$ 可微且满足 $d\alpha/dt=\alpha(t)A$,其中 $A\in\mathcal B(\mathcal H)$ 固定,证明 $\alpha(t)=\alpha(0)e^{tA}$.
\end{enumerate}
\end{exercise}

\begin{exercise}[Fourier 变换实现的酉等价]\label{ex:stone-fourier}
证明 $L^2(\mathbb R^n)$ 上的 $A_j=P_j$、$B_j=-X_j$ 满足指数化对易关系.求出(除相位外唯一的)酉算子 $U:L^2(\mathbb R^n)\to L^2(\mathbb R^n)$,使
\[
Ue^{itA_j}U^{-1}=e^{itX_j},
\qquad
Ue^{itB_j}U^{-1}=e^{itP_j}.
\]
\end{exercise}

\begin{exercise}[Weyl 一参数群]\label{ex:stone-weyl-group}
验证式~\eqref{eq:stone-one-parameter} 中的 $U_{a,b}(t)$ 构成强连续的一参数酉群.
\end{exercise}

\begin{exercise}[有限群上的离散 Stone--von Neumann 定理]\label{ex:finite-stone-von-neumann}
设 $p$ 是素数,$\mathbb Z/p$ 表示模 $p$ 的整数域,$h$ 是其中非零元素.考虑带计数测度的有限维 Hilbert 空间 $L^2(\mathbb Z/p)$.定义调制算子
\[
(Uf)(n)=e^{2\pi in/p}f(n),
\]
以及平移算子
\[
(Vf)(n)=f(n+h).
\]
\begin{enumerate}
\item 验证 $U^p=V^p=I$,并对任意 $l,m\in\mathbb Z$,证明
\[
U^lV^m=e^{-2\pi ilm/p}V^mU^l.
\]
\item 设 $A,B$ 是有限维 Hilbert 空间上的酉算子,满足 $A^p=B^p=I$ 与
\[
A^lB^m=e^{-2\pi ilm/p}B^mA^l.
\]
若只有 $\{0\}$、$\mathcal H$ 在 $A,B$ 下不变,证明存在酉映射 $W:\mathcal H\to L^2(\mathbb Z/p)$,使 $WAW^{-1}=U$、$WBW^{-1}=V$.
\end{enumerate}
提示:若 $v$ 是 $A$ 的特征向量,证明 $B^lv$ 也是特征向量;证明每个特征子空间维数为 $1$,并把相应特征向量识别为 $L^2(\mathbb Z/p)$ 中的 delta 函数.
\end{exercise}

\begin{exercise}[Heisenberg 群与 Weyl 算子]\label{ex:heisenberg-group}
给定 $u\in\mathbb C$ 且 $|u|=1$,以及 $a,b\in\mathbb R^n$,定义 $L^2(\mathbb R^n)$ 上的酉算子
\[
(U_{u,a,b}\psi)(x)=u e^{ia\cdot x}\psi(x+\hbar b).
\]
\begin{enumerate}
\item 验证这类算子在复合运算下构成群,记为 $H_n$.
\item 令 $\widetilde H_n$ 为所有形如
\[
\begin{pmatrix}
1&a_1&\cdots&a_n&c\\
0&1&&0&b_1\\
\vdots&&\ddots&&\vdots\\
0&0&&1&b_n\\
0&0&\cdots&0&1
\end{pmatrix}
\]
的 $(n+2)\times(n+2)$ 矩阵组成的集合,其中 $a_j,b_j,c\in\mathbb R$,除主对角线、第一行和最后一列外的元素为零.验证 $\widetilde H_n$ 在矩阵乘法下构成群,并证明存在满射群同态 $\Phi:\widetilde H_n\to H_n$,其核是离散的.
\end{enumerate}
提示:比较 $H_n$ 与 $\widetilde H_n$ 中的群乘法公式.$\widetilde H_n$ 是 $H_n$ 的万有覆盖群,称为 Heisenberg 群.
\end{exercise}

\begin{exercise}[Segal 平移算子的乘法]\label{ex:segal-translation-product}
直接计算,证明式~\eqref{eq:segal-translation-product} 成立.
\end{exercise}

\begin{exercise}[Segal 平移的强连续性]\label{ex:segal-translation-continuity}
利用支配收敛定理,证明对 $\mathbb C^n$ 上任意全纯多项式 $F$,都有
\[
\lim_{a\to b}\|T_aF-T_bF\|_{L^2(\mathbb C^n,\mu_\hbar)}=0,
\]
其中 $T_a$ 如式~\eqref{eq:segal-translation} 定义.
\end{exercise}

}

\def\ChapterFourteenTheorem{
\section{定理}
\label{sec:stone-theorem}

下面给出 Stone--von Neumann 定理的两个版本:一个适用于满足指数化对易关系的一般算子,另一个适用于不可约作用的特殊情形.

\begin{definition}[不可约作用]\label{def:irreducible-exponentiated-action}
满足指数化对易关系的算子 $A_1,\ldots,A_n,B_1,\ldots,B_n$ 称为在 $\mathcal H$ 上不可约作用,如果唯一同时在所有 $e^{itA_j}$ 与 $e^{itB_j}$ 下不变的闭子空间是 $\{0\}$ 与 $\mathcal H$.
\end{definition}

\begin{proposition}[通常位置与动量的不可约性]\label{prop:usual-position-momentum-irreducible}
通常的位置、动量算子在 $L^2(\mathbb R^n)$ 上不可约作用.
\end{proposition}

该命题的证明延后到本节接近末尾处.

\begin{theorem}[Stone--von Neumann 定理]\label{thm:stone-von-neumann}
设 $A_1,\ldots,A_n,B_1,\ldots,B_n$ 是 $\mathcal H$ 上满足指数化对易关系的自伴算子.则 $\mathcal H$ 可以分解为闭子空间 $\{V_l\}$ 的正交直和,满足:
\begin{enumerate}
\item 对每个 $j,t$,$V_l$ 在 $e^{itA_j}$ 与 $e^{itB_j}$ 下不变;
\item 存在酉算子 $U_l:V_l\to L^2(\mathbb R^n)$,使
\[
U_le^{itA_j}U_l^{-1}=e^{itX_j},
\qquad
U_le^{itB_j}U_l^{-1}=e^{itP_j}.
\]
\end{enumerate}
如果 $A$ 与 $B$ 在 $\mathcal H$ 上不可约作用,则存在单个酉映射 $U:\mathcal H\to L^2(\mathbb R^n)$,使
\[
Ue^{itA_j}U^{-1}=e^{itX_j},
\qquad
Ue^{itB_j}U^{-1}=e^{itP_j}
\]
对所有 $j,t$ 成立.该 $U$ 除乘以绝对值为 $1$ 的常数外是唯一的.
\end{theorem}

上述结果可以用 Heisenberg 群表述;见练习~\ref{ex:heisenberg-group}.

证明定理的策略(正如 von Neumann 1931 年的论文)是沿用第~\ref{sec:stone-heuristic} 节的启发式论证,但把无界升降算子替换为记号~\ref{def:weyl-operator-notation} 中的有界算子 $e^{i(a\cdot A+b\cdot B)}$.取 $\sigma>0$,令
\[
\phi_0(x)=(\pi\sigma)^{-n/4}e^{-|x|^2/(2\sigma)},
\label{eq:stone-ground-state}
\]
它是 $L^2(\mathbb R^n)$ 中的单位向量,可看作频率 $\omega=\hbar/(m\sigma)$ 的 $n$ 维谐振子的基态.对应投影 $|\phi_0\rangle\langle\phi_0|$ 的 Weyl 符号是
\[
f_0(x,p):=Q_{\mathrm W}^{-1}(|\phi_0\rangle\langle\phi_0|)
=2^n e^{-|x|^2/\sigma}e^{-\sigma|p|^2/\hbar^2}.
\label{eq:stone-ground-symbol}
\]
(见练习~\ref{ex:weyl-rank-one-projection}.)

在 $\mathcal H$ 上用 $e^{i(a\cdot A+b\cdot B)}$ 替代 $e^{i(a\cdot X+b\cdot P)}$,可以定义一个广义 Weyl 量子化 $Q$.我们将证明 $\mathcal P:=Q(f_0)$ 是正交投影,并把 $W:=\operatorname{Range}(\mathcal P)$ 作为 $\mathcal H$ 中的基态空间.关键是证明 $\mathcal P$ 非零,甚至其在任意由 $e^{i(a\cdot A+b\cdot B)}$ 不变的非零子空间上的限制都非零.

若 $\{\psi^l\}$ 是 $W$ 的正交归一基,考虑
\[
\psi^l_{a,b}=e^{i(a\cdot A+b\cdot B)}\psi^l.
\]
我们将证明不同 $l$ 的这些向量彼此正交,而固定 $l$ 时,两个向量的内积与 $L^2(\mathbb R^n)$ 中相应向量的内积相同.于是令 $V_l$ 为固定 $l$、让 $a,b$ 变化时 $\psi^l_{a,b}$ 的闭线性张成,就可以构造从 $V_l$ 到 $L^2(\mathbb R^n)$ 的酉映射,并把 $e^{i(a\cdot A+b\cdot B)}$ 交织到通常的 Weyl 算子.所有 $V_l$ 的直和必须是整个 $\mathcal H$:否则其正交补仍由这些算子不变,投影 $\mathcal P$ 在其上的限制非零,从而会产生与所有 $\psi^l$ 正交的 $W$ 中向量,矛盾.

\begin{definition}[广义 Weyl 量子化]\label{def:generalized-weyl-quantization}
设自伴算子 $A_1,\ldots,A_n,B_1,\ldots,B_n$ 在 $\mathcal H$ 上满足指数化对易关系.对任意 $f\in\mathcal S(\mathbb R^{2n})$,定义 $Q(f)\in\mathcal B(\mathcal H)$ 为 Bochner 积分
\[
Q(f)=(2\pi)^{-n}\int_{\mathbb R^{2n}}\widehat f(a,b)e^{i(a\cdot A+b\cdot B)}\,da\,db.
\label{eq:generalized-weyl}
\]
\end{definition}

我们使用 Bochner 积分的标准性质:若连续的 $\mathcal B(\mathcal H)$ 值函数范数可积,则积分有定义;Bochner 积分与施加有界线性变换可交换;并且有 Fubini 定理的相应版本.

\begin{proposition}[广义 Weyl 量子化的基本性质]\label{prop:generalized-weyl-properties}
对任意满足指数化对易关系的算子,定义~\ref{def:generalized-weyl-quantization} 中的 $Q$ 满足:
\begin{enumerate}
\item 若 $f\in\mathcal S(\mathbb R^{2n})$ 为实值,则 $Q(f)$ 自伴.
\item 对任意 $a,b\in\mathbb R^n$ 与 $f\in\mathcal S(\mathbb R^{2n})$,
\[
e^{i(a\cdot A+b\cdot B)}Q(f)=Q(f'),
\qquad
Q(f)e^{i(a\cdot A+b\cdot B)}=Q(f''),
\]
其中
\[
\widehat{f'}(a',b')=e^{i\hbar(a'\cdot b-b'\cdot a)/2}\widehat f(a'-a,b'-b),
\]
\[
\widehat{f''}(a',b')=e^{-i\hbar(a'\cdot b-b'\cdot a)/2}\widehat f(a'-a,b'-b).
\]
\item 对任意 $f,g\in\mathcal S(\mathbb R^{2n})$,
\[
Q(f)Q(g)=Q(f\star g),
\]
其中 $\star$ 是命题~\ref{prop:moyal-product} 中的 Moyal 乘积.
\item 若 $f\in\mathcal S(\mathbb R^{2n})$ 且 $Q(f)=0$,则 $f=0$.
\end{enumerate}
此外,对任意 $a,b\in\mathbb R^n$,有
\[
e^{-i(a\cdot A+b\cdot B)}Q(f)e^{i(a\cdot A+b\cdot B)}=Q(g),
\]
其中
\[
\widehat g(a',b')=e^{i\hbar(a'\cdot b-b'\cdot a)}\widehat f(a',b').
\label{eq:generalized-weyl-conjugation}
\]
\end{proposition}

\begin{proof}
对第 (1) 点,把 $Q(f)$ 写成
\[
(2\pi)^{-n}\int\frac12\left(\widehat f(a,b)e^{i(a\cdot A+b\cdot B)}
+\widehat f(-a,-b)e^{-i(a\cdot A+b\cdot B)}\right)\,da\,db.
\]
第二项作变量代换 $a'=-a,b'=-b$ 后与第一项相同.若 $f$ 实值,$\widehat f(-a,-b)=\overline{\widehat f(a,b)}$,所以积分中的被积算子自伴.

对第 (2) 点,使用式~\eqref{eq:stone-weyl-product}:
\[
e^{i(a\cdot A+b\cdot B)}Q(f)
=(2\pi)^{-n}\int e^{-i\hbar(a\cdot b'-b\cdot a')/2}
\widehat f(a',b')e^{i((a+a')\cdot A+(b+b')\cdot B)},da'\,db'.
\]
令 $a''=a'+a,b''=b'+b$ 并整理,即得到 $f'$ 的公式;第二个公式类似.第 (3) 点与命题~\ref{prop:moyal-product} 的证明完全相同,因为该证明只用到了指数化对易关系.

对第 (4) 点,设 $Q(f)=0$.任意 $\phi,\psi\in\mathcal H$ 与 $a,b\in\mathbb R^n$ 都满足
\[
0=\left\langle e^{i(a\cdot A+b\cdot B)}\phi,
Q(f)e^{i(a\cdot A+b\cdot B)}\psi\right\rangle
=\langle\phi,Q(g)\psi\rangle.
\]
展开 $Q(g)$,得到
\[
0=(2\pi)^{-n}\int e^{i\hbar(a'\cdot b-b'\cdot a)}\widehat f(a',b')
\langle\phi,e^{i(a'\cdot A+b'\cdot B)}\psi\rangle\,da'\,db'.
\label{eq:generalized-weyl-fourier-zero}
\]
这是函数 $\widehat f(a',b')\langle\phi,e^{i(a'\cdot A+b'\cdot B)}\psi\rangle$ 在某点的逆 Fourier 变换.由 Fourier 反演,它对几乎处处的 $(a',b')$ 为零.内积因子关于 $(a',b')$ 连续,并且取 $\phi=e^{i(a_0\cdot A+b_0\cdot B)}\psi$ 可以使它在任意给定点 $(a_0,b_0)$ 非零,因而在该点邻域非零.因此 $\widehat f$ 恒为零,$f$ 也恒为零.最后一个共轭公式由第 (2) 点的两个关系合并得到.
\end{proof}

\begin{lemma}[Gaussian 投影的乘积公式]\label{lem:gaussian-projection}
令
\[
f_0(x,p)=2^n e^{-|x|^2/\sigma}e^{-\sigma|p|^2/\hbar^2},
\qquad \sigma>0.
\]
则对任意 $a,b\in\mathbb R^n$,
\[
Q(f_0)e^{i(a\cdot A+b\cdot B)}Q(f_0)
=e^{-\sigma|a|^2/4}e^{-\hbar^2|b|^2/(4\sigma)}Q(f_0).
\label{eq:gaussian-projection-product}
\]
特别地,$Q(f_0)^2=Q(f_0)$.
\end{lemma}

\begin{proof}
由命题~\ref{prop:generalized-weyl-properties},式~\eqref{eq:gaussian-projection-product} 等价于
\[
f_0\star f'_0=e^{-\sigma|a|^2/4}e^{-\hbar^2|b|^2/(4\sigma)}f_0.
\]
当然可以直接由 $f'_0$ 与 $\star$ 的定义计算;所有积分都是 Gaussian 积分.然而,更简便的办法是只在 $L^2(\mathbb R^n)$ 的通常 Weyl 量子化上验证.由练习~\ref{ex:weyl-rank-one-projection},$Q_{\mathrm W}(f_0)=|\phi_0\rangle\langle\phi_0|$,其中 $\phi_0$ 由式~\eqref{eq:stone-ground-state} 给出.因此
\[
Q_{\mathrm W}(f_0)e^{i(a\cdot X+b\cdot P)}Q_{\mathrm W}(f_0)
=c|\phi_0\rangle\langle\phi_0|,
\qquad
c=\langle\phi_0,e^{i(a\cdot X+b\cdot P)}\phi_0\rangle.
\]
由式~\eqref{eq:weyl-exponential-action},
\[
c=(\pi\sigma)^{-n/2}e^{i\hbar(a\cdot b)/2}
\int_{\mathbb R^n}e^{-|x|^2/(2\sigma)}e^{ia\cdot x}
e^{-|x+\hbar b|^2/(2\sigma)}\,dx.
\]
展开 $|x+\hbar b|^2$,整理指数并使用 Gaussian 积分,得到
\[
c=e^{-\sigma|a|^2/4}e^{-\hbar^2|b|^2/(4\sigma)},
\]
从而证明结论.
\end{proof}

}


现在设 $A,B$ 是满足
\[
[A,B]=i\hbar I
\label{eq:stone-formal-ccr}
\]
的无界自伴算子,并用谱定理定义 $e^{isA},e^{itB}$.形式地把定理~\ref{thm:bch-special} 应用于 $isA,itB$,得到
\[
e^{i(sA+tB)}=e^{ist\hbar/2}e^{isA}e^{itB}
=e^{-ist\hbar/2}e^{itB}e^{isA},
\]
即
\[
e^{isA}e^{itB}=e^{-ist\hbar}e^{itB}e^{isA}.
\label{eq:stone-formal-exponentiated}
\]
必须强调,式~\eqref{eq:stone-formal-exponentiated} 目前只是形式结论,因为它假定有界算子的结果可以转移到无界算子,而这一般是错误的.不过,在“良好”情形中,满足式~\eqref{eq:stone-formal-ccr} 的自伴算子确实可能满足它.

\begin{definition}[指数化对易关系]\label{def:exponentiated-commutation}
设 $A_1,\ldots,A_n$ 与 $B_1,\ldots,B_n$ 是 $\mathcal H$ 上可能无界的自伴算子.如果对任意 $1\leq j,k\leq n$、$s,t\in\mathbb R$,都有
\begin{align*}
e^{isA_j}e^{itA_k}&=e^{itA_k}e^{isA_j},\\
e^{isB_j}e^{itB_k}&=e^{itB_k}e^{isB_j},\\
e^{isA_j}e^{itB_k}&=e^{-ist\hbar\delta_{jk}}e^{itB_k}e^{isA_j},
\end{align*}
则称它们满足指数化对易关系.
\end{definition}

这里的指数由无界自伴算子的谱定理定义,都是定义在整个 $\mathcal H$ 上的酉算子.因此,说指数化对易关系成立,是指这些关系在整个 Hilbert 空间上成立.

\begin{definition}[Weyl 算子记号]\label{def:weyl-operator-notation}
若 $A_1,\ldots,A_n,B_1,\ldots,B_n$ 满足指数化对易关系,对任意 $a,b\in\mathbb R^n$,定义酉算子
\[
e^{i(a\cdot A+b\cdot B)}
:=e^{i\hbar(a\cdot b)/2}e^{ia_1A_1}\cdots e^{ia_nA_n}
e^{ib_1B_1}\cdots e^{ib_nB_n}.
\label{eq:stone-weyl-notation}
\]
\end{definition}

式~\eqref{eq:stone-weyl-notation} 正是形式地把定理~\ref{thm:bch-special} 应用于 $ia\cdot A$、$ib\cdot B$,再形式地使用式~\eqref{eq:exp-commuting} 分解指数所得.通过验证
\[
U_{a,b}(t):=e^{it^2\hbar(a\cdot b)/2}e^{ita_1A_1}\cdots e^{ita_nA_n}
e^{itb_1B_1}\cdots e^{itb_nB_n}
\label{eq:stone-one-parameter}
\]
构成强连续一参数酉群(见练习~\ref{ex:stone-weyl-group}),也可以为这一记号提供严格依据.利用指数化对易关系可直接计算
\[
e^{i(a\cdot A+b\cdot B)}e^{i(a'\cdot A+b'\cdot B)}
=e^{-i\hbar(a\cdot b'-b\cdot a')/2}
e^{i((a+a')\cdot A+(b+b')\cdot B)}.
\label{eq:stone-weyl-product}
\]
特别地,$e^{-i(a\cdot A+b\cdot B)}$ 是 $e^{i(a\cdot A+b\cdot B)}$ 的逆,正如记号所暗示的.

\begin{example}[通常的位置与动量算子]\label{ex:stone-usual-position-momentum}
令 $A_j=X_j$、$B_j=P_j$,其中 $X_j,P_j$ 是 $L^2(\mathbb R^n)$ 上通常的位置、动量算子,则它们满足指数化对易关系.
\end{example}

\begin{proof}
由于 $X_j$ 是乘法算子,$e^{isX_j}$ 是乘以 $e^{isx_j}$ 的算子.另一方面,由例~\ref{ex:translation-unitary-group},
\[
(e^{itP_j}\psi)(x)=\psi(x+t\hbar e_j).
\]
显然 $e^{isX_j}$ 与 $e^{itX_k}$ 对易,$e^{isP_j}$ 与 $e^{itP_k}$ 对易;此外
\[
(e^{itP_k}e^{isX_j}\psi)(x)
=e^{is(x+t\hbar e_k)_j}\psi(x+t\hbar e_k)
=e^{ist\hbar\delta_{jk}}(e^{isX_j}e^{itP_k}\psi)(x),
\]
这正是所需关系.
\end{proof}

\begin{example}[周期边界条件下的反例]\label{ex:stone-wraparound-counterexample}
令 $A$ 为第~\ref{sec:uncertainty-counterexample} 节中的算子,令 $B$ 为该节算子的唯一自伴延拓.则 $A,B$ 不满足指数化对易关系.
\end{example}

\begin{proof}
$A$ 是乘法算子,因此 $e^{isA}$ 是乘以 $e^{isx}$ 的算子.$B=-i\hbar,d/dx$ 带周期边界条件;我们说明 $e^{itB}$ 是“带绕回的平移”.对 $a\in\mathbb R$ 与 $\psi\in L^2([-1,1])$,定义
\[
(S_a\psi)(x)=\psi(x+a-2m_{x,a}),
\]
其中 $m_{x,a}$ 是满足 $-1\leq x+a-2m_{x,a}<1$ 的唯一整数.容易验证 $S_a$ 是 $L^2([-1,1])$ 上的酉算子.

我们声称
\[
e^{itB}=S_{\hbar t}.
\label{eq:stone-wraparound}
\]
事实上,$B$ 的正交归一特征向量基为 $\psi_n(x)=e^{\pi i n x}$($n\in\mathbb Z$),相应特征值为 $\pi n\hbar$.因此谱定理给出 $e^{itB}\psi_n=e^{\pi int\hbar}\psi_n$.另一方面,
\[
(S_a\psi_n)(x)=e^{\pi in(x+a-2m_{x,a})}
=e^{\pi ina}\psi_n(x),
\]
所以 $e^{itB}$ 与 $S_{\hbar t}$ 在每个 $\psi_n$ 上一致,从而在整个 $L^2([-1,1])$ 上一致.

现在考察常函数 $1$:
\[
e^{isA}e^{itB}1=e^{isx},
\]
而
\[
e^{itB}e^{isA}1=e^{is(x+t\hbar-2m_{x,\hbar t})}
=e^{it\hbar s}e^{isx}e^{-2ism_{x,\hbar t}}.
\]
最后一个因子并不总是 $1$,所以指数化对易关系失败.
\end{proof}

\ChapterFourteenTheorem

\section{Segal--Bargmann 空间}
\label{sec:segal-bargmann}

Stone--von Neumann 定理的一个简单例子是 $\mathcal H=L^2(\mathbb R^n)$,取 $A_j=P_j$、$B_j=-X_j$;由练习~\ref{ex:stone-fourier},定理中的酉映射只是定义~\ref{app:def:fourier} 中 Fourier 变换的一个归一化版本.下面构造一个由 $\mathbb C^n$ 上全纯函数组成的更有趣的 Hilbert 空间.

\subsection{升降算子}
\label{subsec:segal-raising-lowering}

若光滑函数 $F:\mathbb C^n\to\mathbb C$ 对每个 $z_j$(固定其他变量)全纯,则称 $F$ 全纯.等价地,若
\[
\frac{\partial F}{\partial\bar z_j}=0,
\qquad
\frac{\partial}{\partial\bar z_j}=\frac12\left(\frac{\partial}{\partial x_j}+i\frac{\partial}{\partial y_j}\right).
\]
算子
\[
\frac{\partial}{\partial z_j}=\frac12\left(\frac{\partial}{\partial x_j}-i\frac{\partial}{\partial y_j}\right)
\]
保持全纯函数空间不变.

Fock 注意到,乘法算子 $z_j$ 与 $\hbar\,\partial/\partial z_j$ 在全纯函数空间上满足
\[
[z_j,z_k]=0,
\qquad
\left[\hbar\frac{\partial}{\partial z_j},\hbar\frac{\partial}{\partial z_k}\right]=0,
\qquad
\left[\hbar\frac{\partial}{\partial z_j},z_k\right]=\hbar\delta_{jk}I.
\label{eq:fock-commutation}
\]
这正是升降算子的基本对易关系;若略去第 11 章式~\eqref{eq:harmonic-ladder-operators} 中的 $\sqrt\hbar$ 归一化,便得到其高维版本.

如果要把 $z_j$ 与 $\hbar\,\partial/\partial z_j$ 理解为互为伴随的升降算子,就需要在全纯函数空间上寻找合适的内积.20 世纪 60 年代初,Segal 与 Bargmann 找到了这样的内积.描述完它之后,我们将把 $z_j$ 与 $\hbar\,\partial/\partial z_j$ 的适当线性组合构造成自伴的位置、动量算子,并验证指数化对易关系与不可约性.

\begin{lemma}[复空间上的分部积分]\label{lem:segal-integration-by-parts}
设 $\mu$ 是 $\mathbb C^n$ 上光滑、严格正的密度,$F,G$ 是足够良好但不必全纯的函数,则
\[
\int_{\mathbb C^n}F(z)\frac{\partial G}{\partial z_j}(z)\mu(z)\,dz
=-\int_{\mathbb C^n}\frac{\partial F}{\partial\bar z_j}(z)G(z)\mu(z)\,dz
-\int_{\mathbb C^n}\frac{\partial\log\mu}{\partial\bar z_j}(z)F(z)G(z)\mu(z)\,dz.
\label{eq:segal-integration-by-parts}
\]
这里 $dz$ 表示 $\mathbb C^n\cong\mathbb R^{2n}$ 上的 Lebesgue 测度.
\end{lemma}

式~\eqref{eq:segal-integration-by-parts} 表明,在 $L^2(\mathbb C^n,\mu)$ 的内积下
\[
\left(\frac{\partial}{\partial z_j}\right)^*
=-\frac{\partial}{\partial\bar z_j}-\frac{\partial\log\mu}{\partial\bar z_j}.
\]
若把伴随限制在全纯函数空间上,第一项因全纯性而消失.

\begin{proof}
将左端积分限制在一个大立方体上,先对 $x_j$ 或 $y_j$ 积分,并分部积分,把 $G$ 上的导数移到 $\overline F\mu$ 上.边界项是在两张相对的面上积分 $F(z)G(z)\mu(z)$;若该函数在无穷远衰减足够快,边界项在极限中消失.又因为
\[
\frac{\partial\overline F}{\partial z_j}=\frac{\partial F}{\partial\bar z_j},
\qquad
\frac{\partial\mu}{\partial z_j}=\frac{\partial\log\mu}{\partial\bar z_j}\mu,
\]
即得式~\eqref{eq:segal-integration-by-parts}.
\end{proof}

我们希望选取密度 $\mu_\hbar$,使 $\partial\log\mu_\hbar/\partial\bar z_j=-z_j/\hbar$.此时 $(\partial/\partial z_j)^*$ 将保持全纯子空间不变,并且在该子空间上等于乘以 $z_j/\hbar$.

\begin{lemma}[Gaussian 密度下的伴随关系]\label{lem:segal-gaussian-adjoint}
令
\[
\mu_\hbar(z)=(\pi\hbar)^{-n}e^{-|z|^2/\hbar}.
\]
若 $F,G$ 是全纯多项式,则
\[
\int_{\mathbb C^n}F(z)\frac{\partial G}{\partial z_j}(z)\mu_\hbar(z)\,dz
=\frac1\hbar\int_{\mathbb C^n}z_jF(z)G(z)\mu_\hbar(z)\,dz.
\label{eq:segal-gaussian-adjoint}
\]
\end{lemma}

\begin{proof}
在式~\eqref{eq:segal-integration-by-parts} 中,全纯性使第一项为零;$\overline F G\mu_\hbar$ 在无穷远快速衰减,边界项消失;最后计算 $\partial\log\mu_\hbar/\partial\bar z_j=-z_j/\hbar$ 即得.
\end{proof}

\begin{definition}[Segal--Bargmann 空间]\label{def:segal-bargmann-space}
记 $\mathcal H L^2(\mathbb C^n,\mu_\hbar)$ 为满足
\[
\|F\|_\hbar=\left(\int_{\mathbb C^n}|F(z)|^2\mu_\hbar(z)\,dz\right)^{1/2}<\infty
\]
的全纯函数 $F$ 的空间.定义升降算子
\[
a_j^*=z_j,
\qquad
a_j=\hbar\frac{\partial}{\partial z_j},
\]
其定义域取为全纯多项式空间.
\end{definition}

由引理~\ref{lem:segal-gaussian-adjoint},对全纯多项式 $F,G$ 有
\[
\langle F,a_jG\rangle_\hbar=\langle a_j^*F,G\rangle_\hbar,
\]
这说明记号 $a_j^*$ 的合理性.该空间也称 Fock 空间;不过在量子场论中,Fock 空间一词还指另一个相关空间,即某个固定 Hilbert 空间上的张量代数的完备化.

\begin{proposition}[Segal--Bargmann 空间的完备性]\label{prop:segal-bargmann-complete}
Segal--Bargmann 空间在范数 $\|\cdot\|_\hbar$ 下完备,关于相应内积构成 Hilbert 空间;全纯多项式在其中稠密.
\end{proposition}

\begin{proof}
对 $z_0\in\mathbb C^n$ 与 $R>0$,令
\[
P_{z_0,R}=\{z\in\mathbb C^n:|z_j-(z_0)_j|<R, j=1,\ldots,n\}.
\]
由幂级数论证,全纯函数 $F$ 在 $z_0$ 的值等于其在多圆盘上的平均值,因而
\[
F(z_0)=\frac1{(\pi R^2)^n}\int_{P_{z_0,R}}F(z)\,dz
=\frac1{(\pi R^2)^n}\int_{P_{z_0,R}}F(z)\frac1{\mu_\hbar(z)}\mu_\hbar(z)\,dz.
\]
Cauchy--Schwarz 不等式给出
\[
|F(z_0)|\leq\frac1{(\pi R^2)^n}
\left(\sup_{z\in P_{z_0,R}}\frac1{\mu_\hbar(z)}\right)^{1/2}
\|F\|_{L^2(\mathbb C^n,\mu_\hbar)}.
\label{eq:segal-point-evaluation}
\]
所以点值映射 $F\mapsto F(z_0)$ 是有界线性泛函.若全纯函数列在 $L^2(\mathbb C^n,\mu_\hbar)$ 中收敛,式~\eqref{eq:segal-point-evaluation} 说明它在紧集上一致收敛,极限仍全纯,因此全纯子空间闭,从而完备.

再取 $F$ 的 Taylor 展开
\[
F(z)=\sum_\nu a_\nu z^\nu,
\label{eq:segal-taylor}
\]
其中 $\nu$ 遍历所有多重指标.利用 Fubini 定理,把每个复变量分别用极坐标积分可知,不同单项式彼此正交.先在逐渐增大的原点多圆盘上积分,再用单调收敛定理,得到
\[
\sum_\nu|a_\nu|^2\|z^\nu\|_{L^2(\mathbb C^n,\mu_\hbar)}^2
=\|F\|_{L^2(\mathbb C^n,\mu_\hbar)}^2<\infty.
\]
因此 Taylor 级数在 $L^2$ 中收敛,而其 $L^2$ 极限与逐点极限同为 $F$,故全纯多项式稠密.
\end{proof}

\ChapterFourteenSegalLater
